% ============================================================================
% WORKSHEET - Section 1.4: Writing Equations for Proportional Relationships
% CCSS 7.RP.A.2c
% ============================================================================
\section{Writing Equations for Proportional Relationships}
\setcounter{wsNumber}{0}

\begin{worksheetSkills}
\begin{itemize}
    \item Write proportional equations in the form $y = kx$
    \item Interpret variables and unit rates in real situations
    \item Use equations to find unknown values efficiently
\end{itemize}
\end{worksheetSkills}

\worksheetHeader{Equation Builder}

\begin{speedDrill}{Write the Rule Fast}
\timerBadge{4}

Write an equation for each proportional relationship.

\bigskip
\renewcommand{\arraystretch}{1.9}
\begin{center}
\begin{tabular}{|C{7cm}|C{4cm}|}
\hline
\textbf{Situation} & \textbf{Equation} \\
\hline
$\$9$ per movie ticket; total cost $c$, tickets $t$ & \answerBlank[3.5cm] \\
\hline
$18$ miles each hour; distance $d$, hours $h$ & \answerBlank[3.5cm] \\
\hline
$\frac{3}{4}$ cup of oats per serving; oats $o$, servings $s$ & \answerBlank[3.5cm] \\
\hline
Table point $(4, 22)$ & \answerBlank[3.5cm] \\
\hline
$2.5$ centimeters on a map for each mile; map length $m$, miles $n$ & \answerBlank[3.5cm] \\
\hline
$12$ points every round; total points $p$, rounds $r$ & \answerBlank[3.5cm] \\
\hline
\end{tabular}
\end{center}
\end{speedDrill}
\worksheetAnswer{(1)~$c = 9t$ because cost equals $9$ dollars for each ticket. (2)~$d = 18h$ because the unit rate is $18$ miles per hour. (3)~$o = \frac{3}{4}s$ because each serving uses $\frac{3}{4}$ cup. (4)~From $(4, 22)$, $k = 22 \div 4 = 5.5$, so the equation is $y = 5.5x$. (5)~$m = 2.5n$ because the map length is $2.5$ centimeters for each mile. (6)~$p = 12r$ because the total points equal $12$ times the number of rounds.}

\worksheetHeader{Equation-to-Situation Match}

\begin{matchIt}{Connect Each Equation to the Correct Story}

\matchRow{$c = 4.25n$}{Each notebook costs $\$4.25$}
\matchRow{$d = 60t$}{A car travels $60$ miles each hour}
\matchRow{$w = 1.5h$}{A faucet fills $1.5$ gallons each hour}
\matchRow{$p = 18m$}{A gym member earns $18$ points each month}
\matchRow{$f = \frac{2}{3}s$}{Each smoothie uses $\frac{2}{3}$ cup of fruit per serving}
\end{matchIt}
\worksheetAnswer{$c = 4.25n$ matches ``Each notebook costs $\$4.25$'' because total cost equals price per item times the number of notebooks. $d = 60t$ matches ``A car travels $60$ miles each hour'' because distance equals rate times time. $w = 1.5h$ matches ``A faucet fills $1.5$ gallons each hour.'' $p = 18m$ matches ``A gym member earns $18$ points each month.'' $f = \frac{2}{3}s$ matches ``Each smoothie uses $\frac{2}{3}$ cup of fruit per serving.'' In every case the coefficient of the variable is the unit rate.}

\worksheetHeader{Fix the Equation Mistakes}

\begin{errorDetective}{Find the Variable Mix-Up}

Marcus wrote an equation for each situation.

\bigskip

Situation: A bike rental costs $\$8$ per hour.
\errorWork{$h = 8c$}
Correct or wrong? \answerBlank[2cm] \quad Correct equation: \answerBlank[5cm]

\bigskip

Situation: A plant grows $3$ centimeters each week. Height is $h$ and weeks is $w$.
\errorWork{$h = w + 3$}
Correct or wrong? \answerBlank[2cm] \quad Correct equation: \answerBlank[5cm]

\bigskip

Situation: $5$ folders cost $\$12.50$.
\errorWork{$c = 5x$}
Correct or wrong? \answerBlank[2cm] \quad Correct equation: \answerBlank[5cm]
\end{errorDetective}
\worksheetAnswer{All three are wrong. (1)~For a bike rental at $\$8$ per hour, the total cost should depend on the hours, so the correct equation is $c = 8h$. Marcus reversed the variables. (2)~A proportional equation uses multiplication, not adding a starting amount. If the plant grows $3$ centimeters each week from zero, then $h = 3w$. (3)~First find the unit rate: $12.50 \div 5 = 2.50$, so the correct equation is $c = 2.5x$ if $x$ is the number of folders. Marcus used the number of folders from one example instead of the unit rate.}

\worksheetHeader{Use the Rule Machine}

\begin{inputOutput}{Complete the Table and Write the Equation}

\textbf{Machine A:} The first pair is $(2, 7)$.

\begin{center}
\begin{tabular}{|C{2.5cm}|C{2.5cm}|}
\hline
\textbf{$x$} & \textbf{$y$} \\
\hline
$2$ & $7$ \\
\hline
$5$ & \answerBlank[2cm] \\
\hline
$8$ & \answerBlank[2cm] \\
\hline
$10$ & \answerBlank[2cm] \\
\hline
\end{tabular}
\end{center}

Equation: \answerBlank[4cm]

\worksheetSep

\textbf{Machine B:} The first pair is $(4, 3)$.

\begin{center}
\begin{tabular}{|C{2.5cm}|C{2.5cm}|}
\hline
\textbf{$x$} & \textbf{$y$} \\
\hline
$4$ & $3$ \\
\hline
$8$ & \answerBlank[2cm] \\
\hline
$12$ & \answerBlank[2cm] \\
\hline
$20$ & \answerBlank[2cm] \\
\hline
\end{tabular}
\end{center}

Equation: \answerBlank[4cm]
\end{inputOutput}
\worksheetAnswer{Machine A: The constant is $k = 7 \div 2 = 3.5$, so the equation is $y = 3.5x$. Then $x = 5$ gives $y = 17.5$, $x = 8$ gives $y = 28$, and $x = 10$ gives $y = 35$. Machine B: The constant is $k = 3 \div 4 = \frac{3}{4}$, so the equation is $y = \frac{3}{4}x$. Then $x = 8$ gives $6$, $x = 12$ gives $9$, and $x = 20$ gives $15$. In each machine, divide $y$ by $x$ first, then multiply the input by that constant.}

\worksheetHeader{Explain the Coefficient}

\begin{explainIt}{What Does the Number in Front Mean?}

A student sees the equation $c = 2.75n$ and says, ``The $2.75$ is just a number we happen to use.''
Explain what the $2.75$ actually means in the situation.

\writeLines{5}

\bigskip

Another student writes $y = 6x$ from the table $(2, 12)$, $(3, 18)$, $(5, 30)$.
Explain how the table proves the equation is correct.

\writeLines{6}
\end{explainIt}
\worksheetAnswer{(1)~In the equation $c = 2.75n$, the number $2.75$ is the constant of proportionality. It means each item costs $\$2.75$, so the total cost is found by multiplying $\$2.75$ by the number of items. The coefficient tells the unit rate, not a random number. (2)~The table proves the equation because the ratio $\frac{y}{x}$ is always $6$: $12 \div 2 = 6$, $18 \div 3 = 6$, and $30 \div 5 = 6$. Since the unit rate is constant, the relationship is proportional with $k = 6$, so the correct equation is $y = 6x$. Each output is exactly six times the input.}

\worksheetDone